%%%
%%% Radical "⿓"
%%%
\section*{Radical 212: ``⿓'' (龙)}\addcontentsline{toc}{section}{Radical 212: ⿓,龙}\addcontentsline{loh}{figure}{\#\#\#\# 212: ⿓}

%%%%%%%%%% 龙 %%%%%%%%%%
\subsection*{龙}\addcontentsline{loh}{figure}{龙}

\begin{Entry}{龙}{5}{⿓}[Kangxi 212]
  \begin{Phonetics}{龙}{long2}[][HSK 3]
    \definition*{s.}{Sobrenome: Long}
    \definition[条]{s.}{dragão; animal mítico e sobrenatural, com chifres, escamas, garras e bigodes, capaz de voar e mergulhar na água, provocar nuvens e chuva | dinossauro; um enorme réptil extinto; referência a certos répteis gigantes da antiguidade | do imperador; dragão como símbolo do imperador; usado na era feudal como símbolo do imperador; também se refere a coisas pertencentes ao imperador | em forma de dragão; com um desenho de dragão; refere"-se a certos objetos que formam uma sequência semelhante a um dragão ou decorados com motivos de dragões}
  \end{Phonetics}
\end{Entry}

\begin{Entry}{龙山}{5,3}{⿓,⼭}
  \begin{Phonetics}{龙山}{long2shan1}
    \definition*{s.}{Condado de Longshan em Xiangxi (湘西) Tujia (土家) e Prefeitura Autônoma de Miao (苗) | Distrito de Longshan da cidade de Liaoyuan (辽源), Jilin (吉林) | Montanha do Dragão}
    \definition{s.}{montanha do dragão}
  \seealsoref{吉林}{ji2lin2}
  \seealsoref{辽源}{liao2yuan2}
  \seealsoref{苗}{miao2}
  \seealsoref{土家}{tu3jia1}
  \seealsoref{湘西}{xiang1xi1}
  \end{Phonetics}
\end{Entry}

\begin{Entry}{龙舟}{5,6}{⿓,⾈}
  \begin{Phonetics}{龙舟}{long2zhou1}[][HSK 7-9]
    \definition{s.}{barco dragão}[龙舟竞渡===corrida de barcos dragão]
  \end{Phonetics}
\end{Entry}

\begin{Entry}{龙虾}{5,9}{⿓,⾍}
  \begin{Phonetics}{龙虾}{long2xia1}
    \definition{s.}{lagosta}
  \end{Phonetics}
\end{Entry}

%%%%% EOF %%%%%

